% LIVE-READY STAGING ONLY, 28 August 2026.
% Release gate: do not integrate this prose until the arbitrary-field
% nonreduced exact-live statement/proof bytes receive exact-live verification.
% The rendered prose below deliberately contains no gate-status language.

\paragraph{Why the source collar has to be identified.}
A tropical chart records the values of finitely many analytic functions; it
does not remember a source point by name.  Two different branches can
therefore have the same image ray, and a remote point can in principle map
into a short segment intended for a local branch calculation.  This matters
for currents.  Euclidean balancing determines the flux carried by a target
ray, whereas the classification argument needs the flux on one specified
source branch.  Equality of ray germs alone does not identify those two
numbers.

There is a second, more elementary, restriction.  A smooth partition
function on the source is not generally an analytic moment coordinate.  On
an annular skeleton the negative logarithm of the norm of an analytic unit
is piecewise integral affine.  A smooth plateau with a genuinely curved
transition cannot have that form.  If a cutoff is retained in a local
calculation, the legal replacement is the factorization in
\eqref{eq:p21p3-gf-source-cutoff-factorization}: one writes
\(\chi=\Phi\circ H_{\mathrm{trop}}\) for an analytic moment block \(H\)
and a smooth function \(\Phi\) on its tropical target.  One appends \(H\)
to the full analytic moment \(F\), pulls superforms back with the notation
\(F^*\beta\), and never declares \(\chi\) itself to be analytic.

\Cref{prop:p21p3-gf-virtual-annulus-bridge} resolves the source ambiguity
in three stages.  At an interior type-\(2\) point \(x\), every branch \(b\)
corresponds to a closed point \(p_b\) of the normal projective integral
residual curve \(C_x\).  A residue function with order one at \(p_b\) lifts
to an analytic unit whose branch-order coordinate is one.  This makes the
outgoing vector primitive over the original field, even when \(p_b\) has
residue degree greater than one.  Weak approximation at finitely many closed
points then supplies a finite analytic block separating every selected ray
from every other active branch.

The source neighborhood is chosen before the final proper localization.  All
virtual-disc components are retained whole and are constant for the first
moment block; only finitely many virtual-annular components can be active.
After those finitely many annular arms have been shortened, the localization
lemma is applied again inside the resulting one-vertex star.  For a selected
branch it produces the full canonical collar specified in
\eqref{eq:p21p3-gf-canonical-full-source-collar} and proves the exact
preimage identity
\eqref{eq:p21p3-gf-source-separated-ray-preimage}.  Consequently a test
supported on the corresponding relative-open ray segment receives source
contribution from precisely that collar.

One may split a virtual annulus after a finite separable extension in order
to compute its slopes, but no current is pulled back to that extension.  If
the split annulus has expansion index \(e_b\), the ordinary current flux and
the calibrated integration-current flux both acquire the same factor
\(e_b\).  Their ratio is therefore defined downstairs over \(K\) and is
independent of the chosen moment presentation.  This is why the canonical
collar, rather than a chosen geometric branch after scalar extension, is the
right source object.
